% Generated draft: requires mathematical and source review.
\documentclass[11pt]{article}
\usepackage[margin=25mm]{geometry}
\usepackage{fontspec}
\setmainfont{Latin Modern Roman}
\usepackage{amsmath,amssymb,mathtools,unicode-math}
\setmathfont{Latin Modern Math}
\usepackage{enumitem,needspace,xurl,hyperref}
\setlength{\parindent}{0pt}
\setlength{\parskip}{5pt}
\setlength{\emergencystretch}{3em}
\newcommand{\N}{\mathbb{N}}
\newcommand{\Z}{\mathbb{Z}}
\newcommand{\Q}{\mathbb{Q}}
\newcommand{\R}{\mathbb{R}}
\newcommand{\C}{\mathbb{C}}
\newcommand{\F}{\mathbb{F}}
\newcommand{\A}{\mathbb{A}}
\newcommand{\PP}{\mathbb{P}}
\newcommand{\E}{\mathbb{E}}
\newcommand{\eps}{\varepsilon}
\newcommand{\dd}{\,\mathrm{d}}
\begin{document}
\section*{UnsolvedMath ID: 10400110}
\textbf{Source problem code:} catalogue entry
\subsection{Definitions and mathematical statement}
Let
\[B_4=\langle\sigma_1,\sigma_2,\sigma_3\mid \sigma_1\sigma_3=\sigma_3\sigma_1,\ \sigma_i\sigma_{i+1}\sigma_i=\sigma_{i+1}\sigma_i\sigma_{i+1}\ (i=1,2)\rangle\]
be the braid group on four strands. With \(t\) an indeterminate and \(R=\Z[t,t^{-1}]\), one standard unreduced Burau convention is
\[\rho_4:B_4\longrightarrow\operatorname{GL}_4(R),\qquad
\rho_4(\sigma_i)=I_{i-1}\oplus\begin{pmatrix}1-t&t\\1&0\end{pmatrix}\oplus I_{3-i}.\]
A representation is faithful if it is injective. The source question is whether
\[\ker\rho_4=\{1\},\quad\text{equivalently}\quad
\forall\beta\in B_4,\ \rho_4(\beta)=I_4\Longrightarrow\beta=1.\]
This is the generic Burau representation, not a representation obtained by assigning a numerical value to \(t\). The customary reduced representation is the action after removing the invariant trivial direction; its generic faithfulness question is equivalent. The source does not choose a matrix convention, so the display fixes a standard representative of that question rather than a specialized parameter.

The supplied research summary reports an affirmative resolution in July 2026 by V. Bharathram, J. S. Birman and T. E. Brendle, with a point-pushing argument and an asserted corollary on faithfulness of the Jones representation of \(B_4\). It also reports faithfulness for \(n\leq3\) attributed to Magnus--Peluso and nonfaithfulness for \(n\geq5\) attributed to Moody, Long--Paton and Bigelow. These are supplied literature assertions; neither the cited 2026 result nor its Jones-representation corollary is independently verified by this statement draft.
\subsection{Short English statement}
Does the generic Burau representation distinguish every nonidentity four-strand braid? The registry marks the question solved and supplies a July 2026 affirmative preprint claim, which remains unverified here.
\subsection{Sources}
\begin{itemize}
\item Local UnsolvedMath registry, permanent ID \(10400110\), exact source code \texttt{catalogue entry}.
\item Ohtsuki, \emph{Problems on invariants of knots and 3-manifolds} (2002), Problem 6.2, PDF page 96: \url{https://msp.org/gtm/2002/04/gtm-2002-04-024s.pdf}. Registry extraction pdf-font-structure; recorded access 2026-07-29.
\item V. Bharathram, J. S. Birman, T. E. Brendle, \emph{The Burau representation of the braid group is faithful for n = 4}, arXiv:2607.05283 (July 2026), as supplied by the research summary.
\item The summary attributes the other-strand results to Magnus--Peluso, Moody, Long--Paton and Bigelow without full citations.
\end{itemize}
\end{document}
